%% ---------------------------------------------------------------------------
\appendix
\section{The Encodings}\label{app:encoding}

The formulas $\MT_k$ and $\G_k$ are built from the same pieces. Let $S$ be
the set of $k$ states and $W=S\setminus\{\Fq\}$.
\begin{enumerate}[(E1)]
\item \emph{Rule table.} A variable $x_{e,s}$ for each neighbourhood $e$ and
each $s\in S$, with the clause $\bigvee_s x_{e,s}$ and the clauses $\neg
x_{e,s}\vee\neg x_{e,s'}$ for $s\ne s'$. The entries $\delta(\Lq,\Lq,\Lq)$ and
$\delta(\Lq,\Lq,\ast)$ are fixed to $\Lq$.
\item \emph{Cells.} A variable $y_{u,s}$ for each encoded cell $u$ and each
$s\in W$, with exactly-one clauses as in (E1). So every encoded cell is in a
working state. The cells at time $0$, the quiescent cells $t<i-1$
(\cref{lem:lightcone}) and the borders are constants, and so is every cell of a
firing row, with the value $\Fq$.
\item \emph{Local rule.} For each encoded cell $u$ with predecessors
$u_1,u_2,u_3$ and each choice $a,b,d$ of their values, let $e=(a,b,d)$ and let
$P$ be the conjunction of $y_{u_1,a}$, $y_{u_2,b}$ and $y_{u_3,d}$, constants
omitted. If $u$ is a working cell, the clauses are $P\wedge x_{e,s}\to y_{u,s}$
for $s\in W$ and $P\to\neg x_{e,\Fq}$; if $u$ is in a firing row, the clause is
$P\to x_{e,\Fq}$.
\item \emph{Return chain.} For an encoded cell $u=(t,i)$ of the line of length
$n$ with $t+i=2n-1$ that is not in its firing row, and the half-line cell
$h=(t,i)$: the clauses $\neg y_{u,s}\vee\neg y_{h,s}$ for $s\in W$.
\item \emph{Pumping.} For $\PUMP(n,n')$: the clause $\bigvee_p z_p$ over the
positions $p$ of $\mathrm{cone}_n(n-1,0)$ (\cref{lem:cone}). The input cells
$v_p$ of $J_n$ and $w_p$ of $J_{n'}$ at $p$ are compared by an auxiliary
variable $z_p$ with the clauses $\neg z_p\vee\neg y_{v_p,s}\vee\neg y_{w_p,s}$
for $s\in W$; constant cells are compared directly.
\end{enumerate}

$\MT_k(\mathcal N;M,N')$ consists of (E1); the half-line cells with $t+i\le M$
and, for each $n\in\mathcal N$, the cells of $R_n$, with (E2) and (E3); (E4) for
the lengths $n$ with $n,n+1\in\mathcal N$; and (E5) for $4\le n<n'\le N'$.
$\G_k(T,N_1,C)$ consists of (E1) and, with (E2)--(E4), the half-line cells with
$t+i\le 2N_1$, the cells of $R_n$ for $2\le n\le T+1$, the cells of $R_n$ with
$\tau\le T$ for $T+2\le n\le N_1$, and the cone of the right end at the firing
time for $T+2\le n\le C$. These regions are closed under predecessors, so with
the table fixed every encoded cell is determined. The clauses (E4) for the
complete lines $2\le n\le T+1$ are not part of \cref{def:G}, but they follow from
(G1) and (G4) by the proof of \cref{lem:return}, in which (G4) gives
$C_\infty(2n-2,1)\ne\Fq$; so the formula is equivalent to \cref{def:G}. For the
formulas $\G_4$ of \cref{thm:four,prop:sharp}, the generator gives one variable
to cells that compute the same function of the table, that is, cells with the
same tree of predecessors down to the same half-line cells and borders; this
does not change the satisfying assignments of the table variables. The formula
of \cref{prop:germs} adds, for each germ $g$, a selector variable $\sigma_g$
with the clauses $\neg\sigma_g\vee x_{e,s}$ for the transitions $(e,s)$ of $g$,
and the clause $\bigvee_g\sigma_g$. The conditions (B1)--(B4) of \cref{prop:balzer} are unit and
binary clauses on the variables $x_{e,s}$.

%% ---------------------------------------------------------------------------
\section{Computational Claims and Artefacts}\label{app:artefacts}

\Cref{tab:artefacts} lists the computational claims of
\cref{sec:applications} with the programs behind them and the published files.
The digests are SHA-256 digests of the uncompressed files, abbreviated to their
first eight hexadecimal digits; the supplementary material lists them in full.

\begin{table}[ht]
\centering\footnotesize
\caption{The computational claims, their programs, and the unsatisfiable
formulas with their LRAT proofs. The programs are in \texttt{src/} and
\texttt{research/} of the supplementary material, and the formulas are the files
\texttt{formulas/*.cnf.xz}; the formulas and the proofs are also archived at
Zenodo (see the data availability statement). In the column Claim, $(\Aq\,\Bq)$ etc.\ is the
involution $I$ of \cref{prop:balzer}. The proof of the last formula was obtained
from a DRAT proof of Kissat (digest \texttt{2dd3e052}, $958$\,MB) with
\textsf{drat-trim}. Each proof is verified by \texttt{lrat-check} and by
\texttt{cake\_lpr} (\texttt{results/cake\_lpr.log}).}
\label{tab:artefacts}
\setlength{\tabcolsep}{3pt}
\begin{tabular}{@{}l>{\raggedright\arraybackslash}p{0.205\textwidth}lrrll@{}}
\toprule
Claim & Programs & Formula & Variables & Clauses & SHA-256 & LRAT proof\\
\midrule
\ref{thm:four}\,(i) & \texttt{gencnf.py} & \texttt{mt4\_9} & $753$ & $17{,}605$ & \texttt{d18d6548} & \texttt{1d511a75}, $177$\,MB\\
\ref{thm:four}\,(ii) & \texttt{gfix.py} & \texttt{k4\_G6} & $919$ & $22{,}960$ & \texttt{0e0ba286} & \texttt{7943a698}, $167$\,MB\\
\ref{thm:four}\,(iii) & \texttt{gfix.py} & \texttt{k4\_G4c} & $874$ & $20{,}461$ & \texttt{124f0954} & \texttt{7f76fe69}, $173$\,MB\\
\ref{prop:germs} & \texttt{germcert.py} & \texttt{k14\_n10} & $17{,}266$ & $559{,}558$ & \texttt{4f9306a8} & \texttt{b2324436}, $11$\,MB\\
\ref{prop:balzer}, $(\Aq\,\Bq)$ & \texttt{balzer.py} & \texttt{balzer\_AB\_10} & $1{,}496$ & $67{,}508$ & \texttt{d87457fe} & \texttt{18809db2}, $380$\,MB\\
\ref{prop:balzer}, $(\Lq\,\Aq)$ & \texttt{balzer.py} & \texttt{balzer\_LA\_9} & $1{,}244$ & $49{,}571$ & \texttt{401ef851} & \texttt{d74de977}, $188$\,MB\\
\ref{prop:balzer}, $(\Lq\,\Bq)$ & \texttt{balzer.py} & \texttt{balzer\_LB\_9} & $1{,}244$ & $49{,}571$ & \texttt{d7ba0d40} & \texttt{5784ee12}, $184$\,MB\\
\ref{prop:balzer}, id & \texttt{balzer\_k.py} & \texttt{balzer\_id\_12} & $2{,}148$ & $114{,}965$ & \texttt{023b99e2} & \texttt{e0231f2f}, $3{,}078$\,MB\\
\midrule
\ref{prop:sharp}\,(a) & \texttt{fsspdfs.c} (enumeration) & \multicolumn{5}{l}{---}\\
\ref{prop:sharp}\,(b) & \texttt{gfix.py}; models checked by \texttt{encoding\_check.py} & \multicolumn{5}{l}{satisfiable formulas}\\
\ref{prop:nearmiss} & \texttt{fsspcheck.c}, \texttt{firstfire.py}, \texttt{failall.py} & \multicolumn{5}{l}{simulations; certificates $(t_0,p)$ in \cref{tab:rules}}\\
\ref{prop:germs} & \texttt{germdfs.c}, \texttt{germdfs\_check.py}, \texttt{germinc.py}, \texttt{germext.c}, \texttt{germ\_k14\_check.py} & \multicolumn{5}{l}{enumeration of the $244$ germs; the $14$ refutations by \cref{lem:left}}\\
\ref{cor:germ12} & \texttt{antidiag.py}, \texttt{germext.c} & \multicolumn{5}{l}{simulation of the cells with $t+i\le1034$}\\
\ref{sec:balzer} & \texttt{balzer\_mazoyer.py}, \texttt{balzer\_rulecheck.py}, \texttt{fsspcheck.c} & \multicolumn{5}{l}{(B1)--(B4) on Mazoyer's solution and on the rules of \cref{tab:balzer}; simulation}\\
\ref{sec:method} & \texttt{encoding\_check.py}, \texttt{encode\_\allowbreak independent.py}, \texttt{validate\_mazoyer.py}, \texttt{gfix.py} & \multicolumn{5}{l}{checks of the encodings, the second encoder; \cref{tab:validation}}\\
\bottomrule
\end{tabular}
\end{table}
